\begin{lexlean}{Proofs.Servable}
\useglossary{lexlean.std.bool@1.2.0}
\useglossary{lexlean.std.nat@1.2.0}
\useglossary{prism.arch@1.0.0}
\importmodule{Kernel}
\title{Text application model}
\begin{semanticmodule}
definition{
  name: breaksServable,
  parameters: [facts : Kernel::Facts],
  result: Bool,
  body: ((((((project facts runner) =? 1) && ((project facts addressed) =? 1)) && (((project facts holds0) =? 1) && ((project facts failed0) =? 0))) && (not (Kernel::eligible0 facts))) || ((((((project facts runner) =? 1) && ((project facts addressed) =? 1)) && (((project facts holds1) =? 1) && ((project facts failed1) =? 0))) && (not (Kernel::eligible1 facts))) || ((((((project facts runner) =? 1) && ((project facts addressed) =? 1)) && (((project facts holds2) =? 1) && ((project facts failed2) =? 0))) && (not (Kernel::eligible2 facts))) || (((((project facts runner) =? 1) && ((project facts addressed) =? 1)) && (((project facts holds3) =? 1) && ((project facts failed3) =? 0))) && (not (Kernel::eligible3 facts))))))
}

theorem{
  name: servableHolds,
  parameters: [facts : Kernel::Facts],
  statement: ((breaksServable facts) = false),
  proof: linear_arithmetic{
    definitions: [
      Kernel::eligible,
      Kernel::eligible0,
      Kernel::eligible1,
      Kernel::eligible2,
      Kernel::eligible3,
      breaksServable
    ]
  },
  axioms: [Classical.choice, Quot.sound, propext]
}
\end{semanticmodule}
\end{lexlean}
